\section{Countermeasures and Discussion}
\label{sec:Countermeasures}

\subsection{Countermeasures}
The core attack of {\name} is that an attacker can alter the output of the magnetic tactile sensor via an external magnetic tag.
To address this issue, we discuss a potential defense strategy grounded in multimodal sensor fusion.

Modern robotic or manipulator systems typically provide additional motion-related signals, such as joint torques, gripper motor currents ~\cite{frankaRobotState}, and the pose, velocity ~\cite{urRTDEGuide}, and motion status of the end-effector ~\cite{robotiq2fManual}. 
Although these signals may not provide the same spatial resolution as tactile sensors, they offer independent sources of physical interaction information. 
Such information can be used to verify whether the tactile response is consistent with the actual motion of the robot.

Specifically, the controller can maintain a short time window for both tactile readings and proprioceptive signals. 
A contact event is accepted only when the increase in tactile force is accompanied by other varying signals, such as an increase in gripper motor current, an increase in joint torque, a decrease in gripper closing velocity, or restricted movement of the end-effector.
If the tactile sensor reports a large normal force while the gripper motor current remains close to the neutral pose, the system can flag this as a potential false contact event and prevent the controller from prematurely stopping the gripper.

Similarly, during a negative normal force injection attack, the tactile sensor may continue to report insufficient contact force even if the gripper has already contacted the object. 
In this case, the system can compare the tactile signal with the gripper motor load, joint torques, and gripper motion status. 
If the gripper motor current or joint torque increases while the gripper closing velocity decreases or the end-effector movement is mechanically restricted, the controller can infer that physical contact has actually occurred despite the low tactile readings. 
Then, the system can stop further increases in force or switch to a safe recovery strategy, rather than continuing to apply pressure on the object.

In practical applications, this defense can be implemented as a lightweight monitoring module. 
For each control cycle, it checks whether the tactile force and other signals show a consistent pattern within a sliding window. 
When inconsistencies persist over several consecutive frames, the robotic system can trigger an alarm, ignore suspicious tactile readings, or enter a conservative fail-safe mode. 
This design aims to reduce false positives while making tactile spoofing attacks more difficult.

\subsection{
Discussion
}


\noindent\textbf{Cross-model applicability.}
As shown in Sec.~\ref {sec:background_limitation} and Sec.~\ref {sec:across_models}, the underlying vulnerability is present across multiple evaluated magnetic tactile sensor models. However, the sensor responses differ across models. Therefore, deterministic force injection remains model-specific and may require offline characterization of the target model and corresponding tag customization. As described in Sec.~\ref {sec:reverse engineering}, the offline characterization can be performed on a separately obtained sensor of the same model and requires no access to the robotic system.


{\emergencystretch=1em
\noindent\textbf{Sensing-modality scope.}
This cross-model generalizability is modality-specific. {\name} relies on magnetic-field manipulation and therefore directly targets magnetic sensing components. It does not directly apply to purely optical, capacitive, or resistive tactile sensors. In multimodal tactile systems, {\name} can affect the magnetic sensing component, while its system-level impact depends on the sensor configuration and fusion mechanism.

\par}
